\documentclass[letterpaper,12pt]{article}
\usepackage[margin=.5in]{geometry}
\usepackage{tikz}
\usepackage{helvet}
\renewcommand{\familydefault}{\sfdefault}
\pagestyle{empty}
\pdfmapfile{+cm.map}
\pdfmapfile{+ps2pk35.map}
\setlength{\parindent}{0pt}
\hyphenpenalty=10000
\exhyphenpenalty=10000
\begin{document}
\null\begin{tikzpicture}[remember picture,overlay,x=1cm,y=1cm]
\begin{scope}[shift={(current page.north west)}]
\node[anchor=north west,inner sep=0pt,font=\fontsize{11.5}{13}\selectfont] at (1.4,-.72) {Week 32 / Squares inside rectangles / Grades 4--5};
\node[anchor=south west,inner sep=0pt,font=\fontsize{9}{11}\selectfont] at (1.4,-27.2) {Bellingham Math Circle / Week 32 / F32-45-v2};
\node[anchor=south east,inner sep=0pt,font=\fontsize{9}{11}\selectfont] at (20.15,-27.2) {1};
\begin{scope}[shift={(1.4,-1.65)}]
\node[anchor=north west,inner sep=0pt,align=left,text width=18.5cm,font=\fontsize{13}{16.38}\selectfont] at (0,0) {Keep square edges on grid lines. Cover rectangles with no gaps, overlaps, or overhang. For the biggest-square rule, cover the largest square at one end, then repeat on the rectangle left over.};
\draw[black,line width=.65pt] (0.5,-2.5) rectangle (4.0,-3.9);
\draw[gray!50,line width=.3pt] (1.2,-2.5) -- (1.2,-3.9);
\draw[gray!50,line width=.3pt] (1.9,-2.5) -- (1.9,-3.9);
\draw[gray!50,line width=.3pt] (2.5999999999999996,-2.5) -- (2.5999999999999996,-3.9);
\draw[gray!50,line width=.3pt] (3.3,-2.5) -- (3.3,-3.9);
\draw[gray!50,line width=.3pt] (0.5,-3.2) -- (4.0,-3.2);
\node[font=\fontsize{12}{14}\selectfont,inner sep=1pt] at (2.25,-2.15) {5};
\node[font=\fontsize{12}{14}\selectfont,inner sep=1pt] at (0.15000000000000002,-3.2) {2};
\draw[->,line width=.8pt] (4.4,-3.2) -- (5.4,-3.2);
\draw[black,line width=.65pt] (6,-2.5) rectangle (9.5,-3.9);
\draw[gray!50,line width=.3pt] (6.7,-2.5) -- (6.7,-3.9);
\draw[gray!50,line width=.3pt] (7.4,-2.5) -- (7.4,-3.9);
\draw[gray!50,line width=.3pt] (8.1,-2.5) -- (8.1,-3.9);
\draw[gray!50,line width=.3pt] (8.8,-2.5) -- (8.8,-3.9);
\draw[gray!50,line width=.3pt] (6,-3.2) -- (9.5,-3.2);
\node[font=\fontsize{12}{14}\selectfont,inner sep=1pt] at (7.75,-2.15) {5};
\node[font=\fontsize{12}{14}\selectfont,inner sep=1pt] at (5.65,-3.2) {2};
\draw[fill=gray!20,line width=1.2pt] (6,-2.5) rectangle (7.4,-3.9);
\node[font=\fontsize{12}{14}\selectfont,inner sep=1pt] at (6.7,-3.2) {2};
\draw[->,line width=.8pt] (10,-3.2) -- (11,-3.2);
\draw[black,line width=.65pt] (11.7,-2.5) rectangle (13.799999999999999,-3.9);
\draw[gray!50,line width=.3pt] (12.399999999999999,-2.5) -- (12.399999999999999,-3.9);
\draw[gray!50,line width=.3pt] (13.1,-2.5) -- (13.1,-3.9);
\draw[gray!50,line width=.3pt] (11.7,-3.2) -- (13.799999999999999,-3.2);
\node[font=\fontsize{12}{14}\selectfont,inner sep=1pt] at (12.75,-2.15) {3};
\node[font=\fontsize{12}{14}\selectfont,inner sep=1pt] at (11.35,-3.2) {2};
\node[anchor=north west,inner sep=0pt,align=left,text width=5cm,font=\fontsize{11}{13.86}\selectfont] at (5.8,-4.5) {cover a 2 by 2 square};
\node[anchor=north west,inner sep=0pt,align=left,text width=5cm,font=\fontsize{11}{13.86}\selectfont] at (11.4,-4.5) {3 by 2 remains};
\node[anchor=north west,inner sep=0pt,align=left,text width=18.6cm,font=\fontsize{14}{17.78}\selectfont] at (0,-6) {\textbf{Problem 1:} Apply the biggest-square rule to both rectangles. Record every square size, including repeated sizes, and compare the last squares.};
\draw[black,line width=.65pt] (1,-8.5) rectangle (11,-14.5);
\draw[gray!50,line width=.3pt] (2,-8.5) -- (2,-14.5);
\draw[gray!50,line width=.3pt] (3,-8.5) -- (3,-14.5);
\draw[gray!50,line width=.3pt] (4,-8.5) -- (4,-14.5);
\draw[gray!50,line width=.3pt] (5,-8.5) -- (5,-14.5);
\draw[gray!50,line width=.3pt] (6,-8.5) -- (6,-14.5);
\draw[gray!50,line width=.3pt] (7,-8.5) -- (7,-14.5);
\draw[gray!50,line width=.3pt] (8,-8.5) -- (8,-14.5);
\draw[gray!50,line width=.3pt] (9,-8.5) -- (9,-14.5);
\draw[gray!50,line width=.3pt] (10,-8.5) -- (10,-14.5);
\draw[gray!50,line width=.3pt] (1,-9.5) -- (11,-9.5);
\draw[gray!50,line width=.3pt] (1,-10.5) -- (11,-10.5);
\draw[gray!50,line width=.3pt] (1,-11.5) -- (11,-11.5);
\draw[gray!50,line width=.3pt] (1,-12.5) -- (11,-12.5);
\draw[gray!50,line width=.3pt] (1,-13.5) -- (11,-13.5);
\node[font=\fontsize{12}{14}\selectfont,inner sep=1pt] at (6.0,-8.15) {10};
\node[font=\fontsize{12}{14}\selectfont,inner sep=1pt] at (0.65,-11.5) {6};
\draw[black,line width=.65pt] (1,-17) rectangle (9,-22);
\draw[gray!50,line width=.3pt] (2,-17) -- (2,-22);
\draw[gray!50,line width=.3pt] (3,-17) -- (3,-22);
\draw[gray!50,line width=.3pt] (4,-17) -- (4,-22);
\draw[gray!50,line width=.3pt] (5,-17) -- (5,-22);
\draw[gray!50,line width=.3pt] (6,-17) -- (6,-22);
\draw[gray!50,line width=.3pt] (7,-17) -- (7,-22);
\draw[gray!50,line width=.3pt] (8,-17) -- (8,-22);
\draw[gray!50,line width=.3pt] (1,-18) -- (9,-18);
\draw[gray!50,line width=.3pt] (1,-19) -- (9,-19);
\draw[gray!50,line width=.3pt] (1,-20) -- (9,-20);
\draw[gray!50,line width=.3pt] (1,-21) -- (9,-21);
\node[font=\fontsize{12}{14}\selectfont,inner sep=1pt] at (5.0,-16.65) {8};
\node[font=\fontsize{12}{14}\selectfont,inner sep=1pt] at (0.65,-19.5) {5};
\draw[gray!45,line width=.35pt] (12,-10) -- (18.5,-10);
\draw[gray!45,line width=.35pt] (12,-12) -- (18.5,-12);
\draw[gray!45,line width=.35pt] (12,-18.5) -- (18.5,-18.5);
\draw[gray!45,line width=.35pt] (12,-20.5) -- (18.5,-20.5);\end{scope}\end{scope}\end{tikzpicture}
\newpage
\null\begin{tikzpicture}[remember picture,overlay,x=1cm,y=1cm]
\begin{scope}[shift={(current page.north west)}]
\node[anchor=north west,inner sep=0pt,font=\fontsize{11.5}{13}\selectfont] at (1.4,-.72) {Week 32 / Squares inside rectangles / Grades 4--5};
\node[anchor=south west,inner sep=0pt,font=\fontsize{9}{11}\selectfont] at (1.4,-27.2) {Bellingham Math Circle / Week 32 / F32-45-v2};
\node[anchor=south east,inner sep=0pt,font=\fontsize{9}{11}\selectfont] at (20.15,-27.2) {2};
\begin{scope}[shift={(1.4,-1.65)}]
\node[anchor=north west,inner sep=0pt,align=left,text width=18.6cm,font=\fontsize{14}{17.78}\selectfont] at (0,0) {\textbf{Problem 2:} Predict the last square in each rectangle, then test. Find a rule that predicts the last square from the two starting side lengths.};
\draw[black,line width=.65pt] (1,-3) rectangle (16,-12);
\draw[gray!50,line width=.3pt] (2,-3) -- (2,-12);
\draw[gray!50,line width=.3pt] (3,-3) -- (3,-12);
\draw[gray!50,line width=.3pt] (4,-3) -- (4,-12);
\draw[gray!50,line width=.3pt] (5,-3) -- (5,-12);
\draw[gray!50,line width=.3pt] (6,-3) -- (6,-12);
\draw[gray!50,line width=.3pt] (7,-3) -- (7,-12);
\draw[gray!50,line width=.3pt] (8,-3) -- (8,-12);
\draw[gray!50,line width=.3pt] (9,-3) -- (9,-12);
\draw[gray!50,line width=.3pt] (10,-3) -- (10,-12);
\draw[gray!50,line width=.3pt] (11,-3) -- (11,-12);
\draw[gray!50,line width=.3pt] (12,-3) -- (12,-12);
\draw[gray!50,line width=.3pt] (13,-3) -- (13,-12);
\draw[gray!50,line width=.3pt] (14,-3) -- (14,-12);
\draw[gray!50,line width=.3pt] (15,-3) -- (15,-12);
\draw[gray!50,line width=.3pt] (1,-4) -- (16,-4);
\draw[gray!50,line width=.3pt] (1,-5) -- (16,-5);
\draw[gray!50,line width=.3pt] (1,-6) -- (16,-6);
\draw[gray!50,line width=.3pt] (1,-7) -- (16,-7);
\draw[gray!50,line width=.3pt] (1,-8) -- (16,-8);
\draw[gray!50,line width=.3pt] (1,-9) -- (16,-9);
\draw[gray!50,line width=.3pt] (1,-10) -- (16,-10);
\draw[gray!50,line width=.3pt] (1,-11) -- (16,-11);
\node[font=\fontsize{12}{14}\selectfont,inner sep=1pt] at (8.5,-2.65) {15};
\node[font=\fontsize{12}{14}\selectfont,inner sep=1pt] at (0.65,-7.5) {9};
\draw[black,line width=.65pt] (1,-15) rectangle (15,-23);
\draw[gray!50,line width=.3pt] (2,-15) -- (2,-23);
\draw[gray!50,line width=.3pt] (3,-15) -- (3,-23);
\draw[gray!50,line width=.3pt] (4,-15) -- (4,-23);
\draw[gray!50,line width=.3pt] (5,-15) -- (5,-23);
\draw[gray!50,line width=.3pt] (6,-15) -- (6,-23);
\draw[gray!50,line width=.3pt] (7,-15) -- (7,-23);
\draw[gray!50,line width=.3pt] (8,-15) -- (8,-23);
\draw[gray!50,line width=.3pt] (9,-15) -- (9,-23);
\draw[gray!50,line width=.3pt] (10,-15) -- (10,-23);
\draw[gray!50,line width=.3pt] (11,-15) -- (11,-23);
\draw[gray!50,line width=.3pt] (12,-15) -- (12,-23);
\draw[gray!50,line width=.3pt] (13,-15) -- (13,-23);
\draw[gray!50,line width=.3pt] (14,-15) -- (14,-23);
\draw[gray!50,line width=.3pt] (1,-16) -- (15,-16);
\draw[gray!50,line width=.3pt] (1,-17) -- (15,-17);
\draw[gray!50,line width=.3pt] (1,-18) -- (15,-18);
\draw[gray!50,line width=.3pt] (1,-19) -- (15,-19);
\draw[gray!50,line width=.3pt] (1,-20) -- (15,-20);
\draw[gray!50,line width=.3pt] (1,-21) -- (15,-21);
\draw[gray!50,line width=.3pt] (1,-22) -- (15,-22);
\node[font=\fontsize{12}{14}\selectfont,inner sep=1pt] at (8.0,-14.65) {14};
\node[font=\fontsize{12}{14}\selectfont,inner sep=1pt] at (0.65,-19.0) {8};\end{scope}\end{scope}\end{tikzpicture}
\newpage
\null\begin{tikzpicture}[remember picture,overlay,x=1cm,y=1cm]
\begin{scope}[shift={(current page.north west)}]
\node[anchor=north west,inner sep=0pt,font=\fontsize{11.5}{13}\selectfont] at (1.4,-.72) {Week 32 / Squares inside rectangles / Grades 4--5};
\node[anchor=south west,inner sep=0pt,font=\fontsize{9}{11}\selectfont] at (1.4,-27.2) {Bellingham Math Circle / Week 32 / F32-45-v2};
\node[anchor=south east,inner sep=0pt,font=\fontsize{9}{11}\selectfont] at (20.15,-27.2) {3};
\begin{scope}[shift={(1.4,-1.65)}]
\node[anchor=north west,inner sep=0pt,align=left,text width=18.6cm,font=\fontsize{14}{17.78}\selectfont] at (0,0) {\textbf{Problem 3:} For each rectangle, find every size of identical whole-grid square that can cover it. Explain why your list is complete.};
\draw[black,line width=.65pt] (1,-3) rectangle (13,-11);
\draw[gray!50,line width=.3pt] (2,-3) -- (2,-11);
\draw[gray!50,line width=.3pt] (3,-3) -- (3,-11);
\draw[gray!50,line width=.3pt] (4,-3) -- (4,-11);
\draw[gray!50,line width=.3pt] (5,-3) -- (5,-11);
\draw[gray!50,line width=.3pt] (6,-3) -- (6,-11);
\draw[gray!50,line width=.3pt] (7,-3) -- (7,-11);
\draw[gray!50,line width=.3pt] (8,-3) -- (8,-11);
\draw[gray!50,line width=.3pt] (9,-3) -- (9,-11);
\draw[gray!50,line width=.3pt] (10,-3) -- (10,-11);
\draw[gray!50,line width=.3pt] (11,-3) -- (11,-11);
\draw[gray!50,line width=.3pt] (12,-3) -- (12,-11);
\draw[gray!50,line width=.3pt] (1,-4) -- (13,-4);
\draw[gray!50,line width=.3pt] (1,-5) -- (13,-5);
\draw[gray!50,line width=.3pt] (1,-6) -- (13,-6);
\draw[gray!50,line width=.3pt] (1,-7) -- (13,-7);
\draw[gray!50,line width=.3pt] (1,-8) -- (13,-8);
\draw[gray!50,line width=.3pt] (1,-9) -- (13,-9);
\draw[gray!50,line width=.3pt] (1,-10) -- (13,-10);
\node[font=\fontsize{12}{14}\selectfont,inner sep=1pt] at (7.0,-2.65) {12};
\node[font=\fontsize{12}{14}\selectfont,inner sep=1pt] at (0.65,-7.0) {8};
\draw[black,line width=.65pt] (1,-14) rectangle (16,-23);
\draw[gray!50,line width=.3pt] (2,-14) -- (2,-23);
\draw[gray!50,line width=.3pt] (3,-14) -- (3,-23);
\draw[gray!50,line width=.3pt] (4,-14) -- (4,-23);
\draw[gray!50,line width=.3pt] (5,-14) -- (5,-23);
\draw[gray!50,line width=.3pt] (6,-14) -- (6,-23);
\draw[gray!50,line width=.3pt] (7,-14) -- (7,-23);
\draw[gray!50,line width=.3pt] (8,-14) -- (8,-23);
\draw[gray!50,line width=.3pt] (9,-14) -- (9,-23);
\draw[gray!50,line width=.3pt] (10,-14) -- (10,-23);
\draw[gray!50,line width=.3pt] (11,-14) -- (11,-23);
\draw[gray!50,line width=.3pt] (12,-14) -- (12,-23);
\draw[gray!50,line width=.3pt] (13,-14) -- (13,-23);
\draw[gray!50,line width=.3pt] (14,-14) -- (14,-23);
\draw[gray!50,line width=.3pt] (15,-14) -- (15,-23);
\draw[gray!50,line width=.3pt] (1,-15) -- (16,-15);
\draw[gray!50,line width=.3pt] (1,-16) -- (16,-16);
\draw[gray!50,line width=.3pt] (1,-17) -- (16,-17);
\draw[gray!50,line width=.3pt] (1,-18) -- (16,-18);
\draw[gray!50,line width=.3pt] (1,-19) -- (16,-19);
\draw[gray!50,line width=.3pt] (1,-20) -- (16,-20);
\draw[gray!50,line width=.3pt] (1,-21) -- (16,-21);
\draw[gray!50,line width=.3pt] (1,-22) -- (16,-22);
\node[font=\fontsize{12}{14}\selectfont,inner sep=1pt] at (8.5,-13.65) {15};
\node[font=\fontsize{12}{14}\selectfont,inner sep=1pt] at (0.65,-18.5) {9};\end{scope}\end{scope}\end{tikzpicture}
\newpage
\null\begin{tikzpicture}[remember picture,overlay,x=1cm,y=1cm]
\begin{scope}[shift={(current page.north west)}]
\node[anchor=north west,inner sep=0pt,font=\fontsize{11.5}{13}\selectfont] at (1.4,-.72) {Week 32 / Squares inside rectangles / Grades 4--5};
\node[anchor=south west,inner sep=0pt,font=\fontsize{9}{11}\selectfont] at (1.4,-27.2) {Bellingham Math Circle / Week 32 / F32-45-v2};
\node[anchor=south east,inner sep=0pt,font=\fontsize{9}{11}\selectfont] at (20.15,-27.2) {4};
\begin{scope}[shift={(1.4,-1.65)}]
\node[anchor=north west,inner sep=0pt,align=left,text width=18.6cm,font=\fontsize{14}{17.78}\selectfont] at (0,0) {\textbf{Problem 4:} A 14 by 8 rectangle loses an 8 by 8 square and leaves a 6 by 8 rectangle. Which whole-number lengths measure both edges exactly before the cut? Which work afterward? Explain why these two lists always match after a biggest-square cut.};
\draw[black,line width=.65pt] (0.5,-4.5) rectangle (9.18,-9.46);
\draw[gray!50,line width=.3pt] (1.12,-4.5) -- (1.12,-9.46);
\draw[gray!50,line width=.3pt] (1.74,-4.5) -- (1.74,-9.46);
\draw[gray!50,line width=.3pt] (2.36,-4.5) -- (2.36,-9.46);
\draw[gray!50,line width=.3pt] (2.98,-4.5) -- (2.98,-9.46);
\draw[gray!50,line width=.3pt] (3.6,-4.5) -- (3.6,-9.46);
\draw[gray!50,line width=.3pt] (4.22,-4.5) -- (4.22,-9.46);
\draw[gray!50,line width=.3pt] (4.84,-4.5) -- (4.84,-9.46);
\draw[gray!50,line width=.3pt] (5.46,-4.5) -- (5.46,-9.46);
\draw[gray!50,line width=.3pt] (6.08,-4.5) -- (6.08,-9.46);
\draw[gray!50,line width=.3pt] (6.7,-4.5) -- (6.7,-9.46);
\draw[gray!50,line width=.3pt] (7.32,-4.5) -- (7.32,-9.46);
\draw[gray!50,line width=.3pt] (7.9399999999999995,-4.5) -- (7.9399999999999995,-9.46);
\draw[gray!50,line width=.3pt] (8.56,-4.5) -- (8.56,-9.46);
\draw[gray!50,line width=.3pt] (0.5,-5.12) -- (9.18,-5.12);
\draw[gray!50,line width=.3pt] (0.5,-5.74) -- (9.18,-5.74);
\draw[gray!50,line width=.3pt] (0.5,-6.359999999999999) -- (9.18,-6.359999999999999);
\draw[gray!50,line width=.3pt] (0.5,-6.98) -- (9.18,-6.98);
\draw[gray!50,line width=.3pt] (0.5,-7.6) -- (9.18,-7.6);
\draw[gray!50,line width=.3pt] (0.5,-8.219999999999999) -- (9.18,-8.219999999999999);
\draw[gray!50,line width=.3pt] (0.5,-8.84) -- (9.18,-8.84);
\node[font=\fontsize{12}{14}\selectfont,inner sep=1pt] at (4.84,-4.15) {14};
\node[font=\fontsize{12}{14}\selectfont,inner sep=1pt] at (0.15000000000000002,-6.98) {8};
\draw[fill=gray!15,line width=1.1pt] (0.5,-4.5) rectangle (5.46,-9.46);
\node[font=\fontsize{13}{15}\selectfont,inner sep=1pt] at (2.98,-6.98) {8 by 8};
\draw[->,line width=.8pt] (9.6,-7) -- (10.7,-7);
\draw[black,line width=.65pt] (11.5,-4.5) rectangle (15.219999999999999,-9.46);
\draw[gray!50,line width=.3pt] (12.12,-4.5) -- (12.12,-9.46);
\draw[gray!50,line width=.3pt] (12.74,-4.5) -- (12.74,-9.46);
\draw[gray!50,line width=.3pt] (13.36,-4.5) -- (13.36,-9.46);
\draw[gray!50,line width=.3pt] (13.98,-4.5) -- (13.98,-9.46);
\draw[gray!50,line width=.3pt] (14.6,-4.5) -- (14.6,-9.46);
\draw[gray!50,line width=.3pt] (11.5,-5.12) -- (15.219999999999999,-5.12);
\draw[gray!50,line width=.3pt] (11.5,-5.74) -- (15.219999999999999,-5.74);
\draw[gray!50,line width=.3pt] (11.5,-6.359999999999999) -- (15.219999999999999,-6.359999999999999);
\draw[gray!50,line width=.3pt] (11.5,-6.98) -- (15.219999999999999,-6.98);
\draw[gray!50,line width=.3pt] (11.5,-7.6) -- (15.219999999999999,-7.6);
\draw[gray!50,line width=.3pt] (11.5,-8.219999999999999) -- (15.219999999999999,-8.219999999999999);
\draw[gray!50,line width=.3pt] (11.5,-8.84) -- (15.219999999999999,-8.84);
\node[font=\fontsize{12}{14}\selectfont,inner sep=1pt] at (13.36,-4.15) {6};
\node[font=\fontsize{12}{14}\selectfont,inner sep=1pt] at (11.15,-6.98) {8};
\draw[gray!45,line width=.35pt] (0,-11.7) -- (18.5,-11.7);
\draw[gray!45,line width=.35pt] (0,-12.85) -- (18.5,-12.85);
\draw[gray!45,line width=.35pt] (0,-14.0) -- (18.5,-14.0);
\node[anchor=north west,inner sep=0pt,align=left,text width=18.6cm,font=\fontsize{14}{17.78}\selectfont] at (0,-16) {\textbf{Problem 5:} Explain why the last square from the biggest-square rule is the largest whole-grid square that can tile the starting rectangle using identical copies.};
\draw[gray!45,line width=.35pt] (0,-19.0) -- (18.5,-19.0);
\draw[gray!45,line width=.35pt] (0,-20.15) -- (18.5,-20.15);
\draw[gray!45,line width=.35pt] (0,-21.3) -- (18.5,-21.3);
\draw[gray!45,line width=.35pt] (0,-22.45) -- (18.5,-22.45);\end{scope}\end{scope}\end{tikzpicture}
\newpage
\null\begin{tikzpicture}[remember picture,overlay,x=1cm,y=1cm]
\begin{scope}[shift={(current page.north west)}]
\node[anchor=north west,inner sep=0pt,font=\fontsize{11.5}{13}\selectfont] at (1.4,-.72) {Week 32 / Squares inside rectangles / Grades 4--5};
\node[anchor=south west,inner sep=0pt,font=\fontsize{9}{11}\selectfont] at (1.4,-27.2) {Bellingham Math Circle / Week 32 / F32-45-v2};
\node[anchor=south east,inner sep=0pt,font=\fontsize{9}{11}\selectfont] at (20.15,-27.2) {5};
\begin{scope}[shift={(1.4,-1.65)}]
\node[anchor=north west,inner sep=0pt,align=left,text width=18.6cm,font=\fontsize{14}{17.78}\selectfont] at (0,0) {\textbf{Problem 6:} Choose a longer side from 9 through 15 and a shorter side from 2 through 8. Which rectangle produces the most different square sizes under the biggest-square rule? Find every rectangle that ties.};
\draw[gray!40,line width=.3pt] (1,-3.35) -- (1,-11.65);
\draw[gray!40,line width=.3pt] (2,-3.35) -- (2,-11.65);
\draw[gray!40,line width=.3pt] (3,-3.35) -- (3,-11.65);
\draw[gray!40,line width=.3pt] (4,-3.35) -- (4,-11.65);
\draw[gray!40,line width=.3pt] (5,-3.35) -- (5,-11.65);
\draw[gray!40,line width=.3pt] (6,-3.35) -- (6,-11.65);
\draw[gray!40,line width=.3pt] (7,-3.35) -- (7,-11.65);
\draw[gray!40,line width=.3pt] (8,-3.35) -- (8,-11.65);
\draw[gray!40,line width=.3pt] (9,-3.35) -- (9,-11.65);
\draw[gray!40,line width=.3pt] (10,-3.35) -- (10,-11.65);
\draw[gray!40,line width=.3pt] (11,-3.35) -- (11,-11.65);
\draw[gray!40,line width=.3pt] (12,-3.35) -- (12,-11.65);
\draw[gray!40,line width=.3pt] (13,-3.35) -- (13,-11.65);
\draw[gray!40,line width=.3pt] (14,-3.35) -- (14,-11.65);
\draw[gray!40,line width=.3pt] (15,-3.35) -- (15,-11.65);
\draw[gray!40,line width=.3pt] (16,-3.35) -- (16,-11.65);
\draw[gray!40,line width=.3pt] (0.85,-3.5) -- (16.15,-3.5);
\draw[gray!40,line width=.3pt] (0.85,-4.5) -- (16.15,-4.5);
\draw[gray!40,line width=.3pt] (0.85,-5.5) -- (16.15,-5.5);
\draw[gray!40,line width=.3pt] (0.85,-6.5) -- (16.15,-6.5);
\draw[gray!40,line width=.3pt] (0.85,-7.5) -- (16.15,-7.5);
\draw[gray!40,line width=.3pt] (0.85,-8.5) -- (16.15,-8.5);
\draw[gray!40,line width=.3pt] (0.85,-9.5) -- (16.15,-9.5);
\draw[gray!40,line width=.3pt] (0.85,-10.5) -- (16.15,-10.5);
\draw[gray!40,line width=.3pt] (0.85,-11.5) -- (16.15,-11.5);
\draw[gray!40,line width=.3pt] (1,-14.35) -- (1,-22.65);
\draw[gray!40,line width=.3pt] (2,-14.35) -- (2,-22.65);
\draw[gray!40,line width=.3pt] (3,-14.35) -- (3,-22.65);
\draw[gray!40,line width=.3pt] (4,-14.35) -- (4,-22.65);
\draw[gray!40,line width=.3pt] (5,-14.35) -- (5,-22.65);
\draw[gray!40,line width=.3pt] (6,-14.35) -- (6,-22.65);
\draw[gray!40,line width=.3pt] (7,-14.35) -- (7,-22.65);
\draw[gray!40,line width=.3pt] (8,-14.35) -- (8,-22.65);
\draw[gray!40,line width=.3pt] (9,-14.35) -- (9,-22.65);
\draw[gray!40,line width=.3pt] (10,-14.35) -- (10,-22.65);
\draw[gray!40,line width=.3pt] (11,-14.35) -- (11,-22.65);
\draw[gray!40,line width=.3pt] (12,-14.35) -- (12,-22.65);
\draw[gray!40,line width=.3pt] (13,-14.35) -- (13,-22.65);
\draw[gray!40,line width=.3pt] (14,-14.35) -- (14,-22.65);
\draw[gray!40,line width=.3pt] (15,-14.35) -- (15,-22.65);
\draw[gray!40,line width=.3pt] (16,-14.35) -- (16,-22.65);
\draw[gray!40,line width=.3pt] (0.85,-14.5) -- (16.15,-14.5);
\draw[gray!40,line width=.3pt] (0.85,-15.5) -- (16.15,-15.5);
\draw[gray!40,line width=.3pt] (0.85,-16.5) -- (16.15,-16.5);
\draw[gray!40,line width=.3pt] (0.85,-17.5) -- (16.15,-17.5);
\draw[gray!40,line width=.3pt] (0.85,-18.5) -- (16.15,-18.5);
\draw[gray!40,line width=.3pt] (0.85,-19.5) -- (16.15,-19.5);
\draw[gray!40,line width=.3pt] (0.85,-20.5) -- (16.15,-20.5);
\draw[gray!40,line width=.3pt] (0.85,-21.5) -- (16.15,-21.5);
\draw[gray!40,line width=.3pt] (0.85,-22.5) -- (16.15,-22.5);\end{scope}\end{scope}\end{tikzpicture}
\end{document}
